% Problem record op_c509b64359fbdc2d. This TeX file is derived from
% database/problems_json/op_c509b64359fbdc2d.json by scripts/sync-tex.mjs; edit the JSON record
% and rerun the script rather than editing this file.
\section{Permanent anticoncentration for complex Gaussian matrices}
\label{sec:op_c509b64359fbdc2d}

\paragraph{Problem.}
Do complex Gaussian permanents satisfy a polynomial lower-tail bound at their root-mean-square scale?
For each integer $n\geq1$, let $X\in\mathbb C^{n\times n}$ have independent entries with density $\pi^{-1}e^{-|z|^2}$.
Write $S_n$ for the permutations of $\{1,\ldots,n\}$ and define
\begin{equation}
 \operatorname{Per}X=\sum_{\pi\in S_n}\prod_{j=1}^nX_{j,\pi(j)}.
 \label{eq:pacc-permanent}
\end{equation}
Does there exist a polynomial $p$, positive on $[1,\infty)^2$, such that the permanent in Eq.~\eqref{eq:pacc-permanent} satisfies
\begin{equation}
 \Pr_X\!\left[|\operatorname{Per}X|<
       \frac{\sqrt{n!}}{p(n,1/\delta)}\right]<\delta
 \qquad(n\geq1,\ 0<\delta<1)?
 \label{eq:pacc-conjecture}
\end{equation}
The same polynomial must work for every pair $(n,\delta)$ in Eq.~\eqref{eq:pacc-conjecture}.

\subsection*{Status}

Solved

\subsection*{Source}

Aaronson and Arkhipov state this as Conjecture~6 in Section~1.2.3 of their arXiv version \sourcecite{ref:pacc-aa}{AA13}.

\subsection*{Progress}

\begin{itemize}
  \item Koehler and Leung's Theorem~1.1 proves that a universal constant $C>0$ satisfies
\begin{equation}
 \sup_{z\in\mathbb C}\Pr\!\left[
 |\operatorname{Per}X-z|\leq t\sqrt{n!}\right]\leq Cnt^2
 \qquad(t>0).
 \label{eq:pacc-small-ball}
\end{equation}
Equation~\eqref{eq:pacc-small-ball} implies Eq.~\eqref{eq:pacc-conjecture}: take $p(n,y)=Kny$ with $K^2>C$.
The resolving result is a July 2026 preprint \sourcecite{ref:pacc-kl}{KL26}.

  \item The earlier Theorem~54 of Aaronson and Arkhipov gives only
$\Pr[|\operatorname{Per}X|^2\geq\theta n!]>(1-\theta)^2/(n+1)$ for $0<\theta<1$.
That moment argument guarantees a fraction of large values; it does not alone control the small-value tail \sourcecite{ref:pacc-aa}{AA13}.
\end{itemize}

\subsection*{References}

\begin{enumerate}[label={},leftmargin=4.8em,labelsep=0.5em]
  \footnotesize
  \item[\textup{[AA13]}]\label{ref:pacc-aa}
  S. Aaronson and A. Arkhipov, "The Computational Complexity of Linear Optics," \emph{Theory of Computing} \textbf{9}(4), 143--252 (2013). \href{https://doi.org/10.4086/toc.2013.v009a004}{doi:10.4086/toc.2013.v009a004}; \href{https://arxiv.org/abs/1011.3245}{arXiv:1011.3245}.

  \item[\textup{[KL26]}]\label{ref:pacc-kl}
  F. Koehler and P. K. Leung, "Anticoncentration of the Permanent in Ginibre Ensembles," arXiv preprint (2026). \href{https://arxiv.org/abs/2607.20329v1}{arXiv:2607.20329v1}.
\end{enumerate}

\subsection*{Comment}

The archived anticoncentration question is resolved by a preprint; no journal publication is asserted.
It is distinct from the Permanent-of-Gaussians hardness conjecture.
Combining this result with Aaronson--Arkhipov Theorem~7 makes their additive squared-permanent and multiplicative permanent estimation tasks polynomial-time equivalent \sourcecite{ref:pacc-aa}{AA13}.
It does not prove either task hard.

\subsection*{Field}

Quantum algorithm

\subsection*{Topic}

Boson sampling; Computational complexity and computability

\subsection*{ID}

\texttt{op\_c509b64359fbdc2d}
